% GENERATED FILE. Edit canonical lessons, then run npm run generate:books.
% canonical-source-hash: fnv1a64:a5a7b80fec76c2e0
% canonical-lessons: GU-C181-utaval-karvi, GU-C181-multavi-rakhvu, GU-C181-talvu-avoid, GU-C181-rad-karvu, GU-C181-lambavvu

\chapter{Hurry and delay}
\label{ch:gu-hurry-and-delay}
\hlchaptermodality{181}

\begin{chapteropening}
\textbf{By the end of this chapter:} \emph{I can say that plans will be hurried, put off, cancelled or extended.}

Complete the last lesson of chapter 181: I can say that plans will be hurried, put off, cancelled or extended.
\end{chapteropening}

\section[\texorpdfstring{\gu{ઉતાવળ} \gu{કરવી} (utāvaḷ karvī)}{ઉતાવળ કરવી (utāvaḷ karvī)}]{\gu{ઉતાવળ} \gu{કરવી} (utāvaḷ karvī) — to hurry}

\label{lesson:GU-C181-utaval-karvi}



\begin{quote}
Before the new one: say the Gujarati for to stay on, then the Gujarati for to pass.
\end{quote}

\subsection*{You'll want to know: \gu{ઉતાવળ} \gu{કરવી}}
\textbf{\gu{ઉતાવળ} \gu{કરવી}} — \emph{utāvaḷ karvī} — \textquotedblleft{}to hurry\textquotedblright{}.

\begin{quote}
\textbf{\gu{આજે} \gu{તે} \gu{ઉતાવળ} \gu{કરશે}.} — \emph{āje te utāvaḷ karśe} — \textquotedblleft{}Today she will be in a hurry.\textquotedblright{}
\end{quote}

\begin{cousinweb}
\textbf{\gu{ઉતાવળ} \gu{કરવી}} (\emph{utāvaḷ karvī}) is built from the feminine noun \emph{utāvaḷ}, \textquotedblleft{}haste\textquotedblright{}, and \emph{karvī}, the feminine of \emph{karvũ}, \textquotedblleft{}to do\textquotedblright{}. CDIAL lists \emph{utāvaḷ}, glossed \textquotedblleft{}impetuous\textquotedblright{}, under \emph{*uttāpala}, a reconstructed form on the root \emph{tap}, \textquotedblleft{}to heat\textquotedblright{} (CDIAL 1788).
\end{cousinweb}

\begin{grammarlens}[title={What you've built}]
One more word for everyday life at A2.
\end{grammarlens}

\subsection*{Your turn}
\begin{itemize}
\raggedright
  \item \emph{Say it:} \emph{utāvaḷ karvī}
  \item \emph{Say it:} \emph{utāvaḷ karvī}, once more
  \item \emph{Say it:} \emph{āje te utāvaḷ karśe} — Today she will be in a hurry.
  \item \emph{Recall:} say the Gujarati for to stay on, then the Gujarati for to pass, then say \emph{utāvaḷ karvī} again
\end{itemize}

\begin{tcolorbox}[breakable,colback=teal!4,colframe=teal!35!black,title={Before you move on}]
What does \gu{ઉતાવળ} \gu{કરવી} mean? (\textquotedblleft{}To hurry\textquotedblright{}.) Say it once more.
\end{tcolorbox}
\section[\texorpdfstring{\gu{મુલતવી} \gu{રાખવું} (multavī rākhvũ)}{મુલતવી રાખવું (multavī rākhvũ)}]{\gu{મુલતવી} \gu{રાખવું} (multavī rākhvũ) — to postpone}

\label{lesson:GU-C181-multavi-rakhvu}



\begin{quote}
Before the new one: say the Gujarati for to pass, then the Gujarati for to hurry.
\end{quote}

\subsection*{You'll want to know: \gu{મુલતવી} \gu{રાખવું}}
\textbf{\gu{મુલતવી} \gu{રાખવું}} — \emph{multavī rākhvũ} — \textquotedblleft{}to postpone\textquotedblright{}.

\begin{quote}
\textbf{\gu{અમે} \gu{મુલાકાત} \gu{મુલતવી} \gu{રાખીશું}.} — \emph{ame mulākāt multavī rākhīśũ} — \textquotedblleft{}We will postpone the meeting.\textquotedblright{}
\end{quote}

\begin{cousinweb}
\textbf{\gu{મુલતવી} \gu{રાખવું}} (\emph{multavī rākhvũ}) is built from \emph{multavī}, \textquotedblleft{}postponed\textquotedblright{}, and the inherited \emph{rākhvũ}, \textquotedblleft{}to keep\textquotedblright{}. \emph{Multavī} is a loan through Persian, usually traced to Arabic \emph{multawī}, \textquotedblleft{}twisted\textquotedblright{}, on the root \emph{l-w-y}, \textquotedblleft{}to twist; to wait for\textquotedblright{} (Lane); Lane does not give the sense \textquotedblleft{}postponed\textquotedblright{}.
\end{cousinweb}

\begin{grammarlens}[title={What you've built}]
One more word for everyday life at A2.
\end{grammarlens}

\subsection*{Your turn}
\begin{itemize}
\raggedright
  \item \emph{Say it:} \emph{multavī rākhvũ}
  \item \emph{Say it:} \emph{multavī rākhvũ}, once more
  \item \emph{Say it:} \emph{ame mulākāt multavī rākhīśũ} — We will postpone the meeting.
  \item \emph{Recall:} say the Gujarati for to pass, then the Gujarati for to hurry, then say \emph{multavī rākhvũ} again
\end{itemize}

\begin{tcolorbox}[breakable,colback=teal!4,colframe=teal!35!black,title={Before you move on}]
What does \gu{મુલતવી} \gu{રાખવું} mean? (\textquotedblleft{}To postpone\textquotedblright{}.) Say it once more.
\end{tcolorbox}
\section[\texorpdfstring{\gu{ટાળવું} (ṭāḷvũ)}{ટાળવું (ṭāḷvũ)}]{\gu{ટાળવું} (ṭāḷvũ) — to avoid, to put off}

\label{lesson:GU-C181-talvu-avoid}



\begin{quote}
Before the new one: say the Gujarati for to hurry, then the Gujarati for to postpone.
\end{quote}

\subsection*{You'll want to know: \gu{ટાળવું}}
\textbf{\gu{ટાળવું}} — \emph{ṭāḷvũ} — \textquotedblleft{}to avoid, to put off\textquotedblright{}.

\begin{quote}
\textbf{\gu{તે} \gu{ફરી} \gu{કામ} \gu{ટાળશે}.} — \emph{te pharī kām ṭāḷśe} — \textquotedblleft{}He will dodge the work again.\textquotedblright{}
\end{quote}

\begin{cousinweb}
\textbf{\gu{ટાળવું}} (\emph{ṭāḷvũ}) is a tadbhava. It comes from Sanskrit \emph{ṭālayati}, the causative of \emph{ṭalati}, \textquotedblleft{}moves aside\textquotedblright{}, and CDIAL glosses the Gujarati verb \textquotedblleft{}to avoid\textquotedblright{} (CDIAL 5450).
\end{cousinweb}

\begin{grammarlens}[title={What you've built}]
One more word for everyday life at A2.
\end{grammarlens}

\subsection*{Your turn}
\begin{itemize}
\raggedright
  \item \emph{Say it:} \emph{ṭāḷvũ}
  \item \emph{Say it:} \emph{ṭāḷvũ}, once more
  \item \emph{Say it:} \emph{te pharī kām ṭāḷśe} — He will dodge the work again.
  \item \emph{Recall:} say the Gujarati for to hurry, then the Gujarati for to postpone, then say \emph{ṭāḷvũ} again
\end{itemize}

\begin{tcolorbox}[breakable,colback=teal!4,colframe=teal!35!black,title={Before you move on}]
What does \gu{ટાળવું} mean? (\textquotedblleft{}To avoid, to put off\textquotedblright{}.) Say it once more.
\end{tcolorbox}
\section[\texorpdfstring{\gu{રદ} \gu{કરવું} (rad karvũ)}{રદ કરવું (rad karvũ)}]{\gu{રદ} \gu{કરવું} (rad karvũ) — to cancel}

\label{lesson:GU-C181-rad-karvu}



\begin{quote}
Before the new one: say the Gujarati for to postpone, then the Gujarati for to avoid.
\end{quote}

\subsection*{You'll want to know: \gu{રદ} \gu{કરવું}}
\textbf{\gu{રદ} \gu{કરવું}} — \emph{rad karvũ} — \textquotedblleft{}to cancel\textquotedblright{}.

\begin{quote}
\textbf{\gu{વરસાદમાં} \gu{તેઓ} \gu{કાર્યક્રમ} \gu{રદ} \gu{કરશે}.} — \emph{varsādmā̃ teo kāryakram rad karśe} — \textquotedblleft{}In the rain they will cancel the programme.\textquotedblright{}
\end{quote}

\begin{cousinweb}
\textbf{\gu{રદ} \gu{કરવું}} (\emph{rad karvũ}) is built from \emph{rad}, \textquotedblleft{}cancelled\textquotedblright{}, and the inherited \emph{karvũ}, \textquotedblleft{}to do\textquotedblright{}. \emph{Rad} is a loan from Arabic \emph{radd}, \textquotedblleft{}a turning back, a rejecting\textquotedblright{}, through Persian; its root is \emph{r-d-d}, \textquotedblleft{}to send back\textquotedblright{} (Lane).
\end{cousinweb}

\begin{grammarlens}[title={What you've built}]
One more word for everyday life at A2.
\end{grammarlens}

\subsection*{Your turn}
\begin{itemize}
\raggedright
  \item \emph{Say it:} \emph{rad karvũ}
  \item \emph{Say it:} \emph{rad karvũ}, once more
  \item \emph{Say it:} \emph{varsādmā̃ teo kāryakram rad karśe} — In the rain they will cancel the programme.
  \item \emph{Recall:} say the Gujarati for to postpone, then the Gujarati for to avoid, then say \emph{rad karvũ} again
\end{itemize}

\begin{tcolorbox}[breakable,colback=teal!4,colframe=teal!35!black,title={Before you move on}]
What does \gu{રદ} \gu{કરવું} mean? (\textquotedblleft{}To cancel\textquotedblright{}.) Say it once more.
\end{tcolorbox}
\section[\texorpdfstring{\gu{લંબાવવું} (lambāvvũ)}{લંબાવવું (lambāvvũ)}]{\gu{લંબાવવું} (lambāvvũ) — to extend, to lengthen}

\label{lesson:GU-C181-lambavvu}



\begin{quote}
Before the new one: say the Gujarati for to avoid, then the Gujarati for to cancel.
\end{quote}

\subsection*{You'll want to know: \gu{લંબાવવું}}
\textbf{\gu{લંબાવવું}} — \emph{lambāvvũ} — \textquotedblleft{}to extend, to lengthen\textquotedblright{}.

\begin{quote}
\textbf{\gu{અમે} \gu{રજા} \gu{બે} \gu{દિવસ} \gu{લંબાવીશું}.} — \emph{ame rajā be divas lambāvīśũ} — \textquotedblleft{}We will extend the holiday by two days.\textquotedblright{}
\end{quote}

Built on \textbf{\gu{લાંબું}}, \emph{lā̃bũ}, long.

\begin{cousinweb}
\textbf{\gu{લંબાવવું}} (\emph{lambāvvũ}) is built on the inherited \emph{lā̃bũ}, \textquotedblleft{}long\textquotedblright{}, with the causative \emph{-āv-}. CDIAL traces \emph{lā̃bũ} to Sanskrit \emph{lamba}, \textquotedblleft{}pendent\textquotedblright{} (CDIAL 10951).
\end{cousinweb}

\begin{grammarlens}[title={What you've built}]
One more word for everyday life at A2.
\end{grammarlens}

\subsection*{Your turn}
\begin{itemize}
\raggedright
  \item \emph{Say it:} \emph{lambāvvũ}
  \item \emph{Say it:} \emph{lambāvvũ}, once more
  \item \emph{Say it:} \emph{ame rajā be divas lambāvīśũ} — We will extend the holiday by two days.
  \item \emph{Recall:} say the Gujarati for to avoid, then the Gujarati for to cancel, then say \emph{lambāvvũ} again
\end{itemize}

\begin{tcolorbox}[breakable,colback=teal!4,colframe=teal!35!black,title={Before you move on}]
What does \gu{લંબાવવું} mean? (\textquotedblleft{}To extend, to lengthen\textquotedblright{}.) Say it once more.
\end{tcolorbox}
